%%%%%%%%%%%%%%%%%%%%%%%%%%%%%%%%%%%%%%%%%%%%%%%%%%%%%%%%%%%
%%%%%%%%%%%%%%%%%%%%%%%%%%%%%%%%%%%%%%%%%%%%%%%%%%%%%%%%%%%
%%%%%%%%%%%%%%%%%%%%%%%%%%%%%%%%%%%%%%%%%%%%%%%%%%%%%%%%%%%
% Consistency of ITR estimator
\newcommand{\ITRconsistency}{
\begin{theorem} (Consistency of the individualized treatment regime estimator) 
\label{thm:ITR_Consistency}
Assume Assumptions \ref{assump:CuiAssump1}-\ref{assump:CuiAssump2} hold. Let $\doptHAT$ denote the optimal ITR estimator that is built following Algorithm \ref{alg:Overview} and Assumptions \ref{assump:boundeddensity}-\ref{assump:boundedV0}. Assume Assumptions \ref{assump:Causal1}-\ref{assump:Causal3} hold and that the random forests are built for each treatment arm. Then, for %each $\phi = \phi^{\mu}, \phi^{\sigma,t_0}$ where 
$\phi^\mu = \begin{pmatrix}
    \phi^\mu_1 \\ \phi^\mu_2
\end{pmatrix},$ % and $\phi^{\sigma, t_0} = \begin{pmatrix} \phi^{\sigma, t_0}_1 \\ \phi^{\sigma, t_0}_2 \end{pmatrix},$
\[ \lim_{n \to \infty} \| \mathbb{V}(\doptTRUE(Z)) - \mathbb{V}(\doptHAT(Z)) \| \xrightarrow{p} 0\]
where $\|\cdot\|$ is defined as the $L_2$ norm.
\end{theorem}
}

\newcommand{\ITRconsistencyPROOF}{
\begin{proof}[Proof of Theorem \ifnum\jasa=1 1\else \ref{thm:ITR_Consistency}\fi]
%We know that, by properties of norm, $\left\| \begin{pmatrix} a_1 \\ a_2 \end{pmatrix} + \begin{pmatrix} b_1 \\ b_2 \end{pmatrix}\right\| \leq \left\| \begin{pmatrix} a_1 \\ a_2 \end{pmatrix} \right\| + \left\| \begin{pmatrix} b_1 \\ b_2 \end{pmatrix} \right\|$ and $\left\| c\begin{pmatrix} a_1 \\ a_2 \end{pmatrix} \right\| =  |c| \left\|\begin{pmatrix} a_1 \\ a_2 \end{pmatrix}\right\|$ where $c$ is a scalar. Also, by property of expectation, we know that $\|\mathbb{E}(\cdot)\| \leq \mathbb{E}(\|\cdot\|)$.

Note that
\begin{equation*}
\begin{split}
\left\| \mathbb{V}(\doptTRUE(Z)) - \mathbb{V}(\doptHAT(Z)) \right\| 
% line 1
&= \left\|\sum_{a \in \mathcal{A}} \mathbb{E}_Z\left[ \valueTRUE
\cdot \left( \mathbf{1}\{\doptTRUE(Z) = a\} - \mathbf{1}\{\doptHAT(Z) = a\} \right)
\right] \right\| \\
% line 2
&\leq \sum_{a \in \mathcal{A}} \left\|\mathbb{E}_Z\left[ \valueTRUE
\cdot \left( \mathbf{1}\{\doptTRUE(Z) = a\} - \mathbf{1}\{\doptHAT(Z) = a\} \right)
\right] \right\| \\
% line 3
&\leq \sum_{a \in \mathcal{A}} \mathbb{E}_Z \left\| \left[ \valueTRUE
\cdot \left( \mathbf{1}\{\doptTRUE(Z) = a\} - \mathbf{1}\{\doptHAT(Z) = a\} \right)
\right] \right\| \\
% line 4
&\leq \sum_{a \in \mathcal{A}} \mathbb{E}_Z \left[ \left\| \valueTRUE \right\|
\cdot \left| \mathbf{1}\{\doptTRUE(Z) = a\} - \mathbf{1}\{\doptHAT(Z) = a\} \right|
\right] \\
% line 5
&= \mathbb{E}_Z \left[ \sum_{a \in \mathcal{A}} \left( \mathbf{1} \{\doptTRUE(Z) \neq \doptHAT(Z)\} \left\|\valueTRUE\right\| \right) \right] \\
% line 6
&\leq \mathbb{E}_Z \left( \mathbf{1} \{\doptTRUE(Z) \neq \doptHAT(Z)\} \right) \sup_{a,Z} \left\| \valueTRUE \right\| \\
% line 7
&= \text{Pr}(\doptTRUE(Z) \neq \doptHAT(Z) | \doptHAT) \sup_{a,Z} \left\| \valueTRUE \right\|\\
% line 8
&\leq \sevenbound \text{Pr}(\doptTRUE(Z) \neq \doptHAT(Z) | \doptHAT).\\
\end{split} % split gives one number; align gives a number per line
\end{equation*} The third to last line holds because 1) only one term in the summation will have a non-zero contribution and 2) the supremum represents the ``worst case scenario" over all fixed combinations of $a$ and $z$ and thus does not depend on the random variables. The last line holds because we assumed $\sup_{a,Z} \left\| \valueTRUE \right\|$ is bounded by $\sevenbound$ via Assumption \ifnum\jasa=1 11\else \ref{assump:boundedV0} \fi. Thus, the norm of this distance, $\mathbb{V}(\doptTRUE) - \mathbb{V}(\doptHAT)$, goes to $0$ in probability if $\text{Pr}(\doptTRUE(Z) \neq \doptHAT(Z)|\hat{d}_\text{opt}) \xrightarrow{p} 0$.\\

We complete the proof by first showing
\begin{align}
    \label{eqn:step1}
        \sum_{a \in \mathcal{A}} \left| I\{ \valueoneTRUE - (1-\tol)\valueonemaxTRUE \ge 0 \} - I\{ \valueoneHAT - (1-\tol)\valueonemaxHAT \ge 0 \} \right| \rightarrow 0.
        \end{align}
    %(Note: we take the difference because both are random (one relies on Z and the other is an estimator)).
The proof of \cref{eqn:step1} will be presented shortly. Now, \cref{eqn:step1} implies 
\begin{align}
    \label{eqn:star}
        \argmax_{a \in \mathcal{A}: \valueoneHAT \ge (1-\tol)\valueonemaxHAT} \valuetwoHAT \xrightarrow{p} \argmax_{a \in \mathcal{A}: \valueoneTRUE \ge (1-\tol)\valueonemaxTRUE} \valuetwoHAT.
    \end{align} 
Theorem \ifnum \jasa=1 3 \else \ref{thm:RCIF_Consistency} \fi in our paper states that RCIFs are consistent. %replace <- with theorem when typed up.
Thus, by the Continuous Mapping Theorem, $\valuetwoHAT \xrightarrow{p} \valuetwoTRUE$. %$\valuetwoHAT$ is bounded.
Then,
\begin{align}
\label{eqn:step2}
        \sup_{a \in \mathcal{A}} \, \mathbb{E} \, \left| \valuetwoHAT - \valuetwoTRUE \right| \xrightarrow{p} 0,
    \end{align} 
which implies that the right hand side of \cref{eqn:star} converges to \\$\argmax_{a \in \mathcal{A}: \valueoneTRUE \ge (1-\tol)\valueonemaxTRUE} \valuetwoTRUE$. Thus, since $Z$ is random, $$\text{Pr}(\doptTRUE(Z) \neq \doptHAT(Z) | \doptHAT) \xrightarrow{p} 0.$$ This completes the proof of Theorem \ifnum\jasa=1 1\else \ref{thm:ITR_Consistency}\fi.
\end{proof}
}

\newcommand{\eqnsteponePROOF}{
%\textbf{Proof of \cref{eqn:step1}:}
\begin{proof}[Proof of \cref{eqn:step1}]
Let $A_0 = \valueoneTRUE - (1-\tol)\valueonemaxTRUE$
and $\widehat{A}_n = \valueoneHAT - (1-\tol)\valueonemaxHAT$. These are functions of random variable $Z$. Then, we can re-write \cref{eqn:step1} as 
$$\sum_{a \in \mathcal{A}} \left| I\{ A_0 \ge 0 \} - I \{ \widehat{A}_n \ge 0 \} \right| \rightarrow 0.$$
\begin{equation}
\label{eqn:terms12}
\begin{split}
    \left| I\{ \widehat{A}_n \ge 0 \} - I \{A_0 \ge 0 \} \right| 
    &= I\{ \widehat{A}_n < 0 \le A_0 \} + I \{A_0 < 0 \le \widehat{A}_n \} \\
    &= I\{ \widehat{A}_n - A_0 < -A_0 \leq 0 \} + I \{0 < -A_0 \le \widehat{A}_n - A_0 \} \\
    &= I\{ 0 \le A_0 < -(\widehat{A}_n - A_0) \} + I \{- (\widehat{A}_n - A_0) \le A_0 < 0 \}. \\
\end{split}
\end{equation}
In \cref{eqn:terms12}, the first equality holds because these are the only possible indicators such that their sum is equal to $1$. We get the 2nd equality by subtracting $A_0$ on both sides within the indicator functions. The third equality is due to multiplying by $-1$ and flipping the inequalities within the indicators.

Since we know $\valueoneHAT \xrightarrow{p} \valueoneTRUE$ and $\valueonemaxHAT \xrightarrow{p} \valueonemaxTRUE$ as shown in \cite{cho_multi-stage_2023}; and $\widehat{A}_n$ and $A_0$ are linear combinations of $\valueoneHAT$ and $\valueonemaxHAT$, and $\valueoneTRUE$ and $\valueonemaxTRUE$, respectively, then %Theorem \ref{thm:consistency_lincomb},
$\widehat{A}_n \xrightarrow{p} A_0$ because the linear combination of consistent estimators is consistent. This means
\begin{align}
\label{eqn:AnA0} 
\widehat{A}_n - A_0 \xrightarrow{p} 0.
\end{align}

By definition of convergence in probability, \cref{eqn:AnA0} implies $\exists\epsilon_n: Pr(|\widehat{A}_n - A_0|>\epsilon_n) \rightarrow 0 \text{ as } \epsilon_n \rightarrow 0$ for a fixed sequence $\epsilon_n$. Then, using Assumption \ifnum\jasa=1 10 \else \ref{assump:boundeddensity} \fi and \cref{eqn:AnA0}, we can bound \cref{eqn:terms12}. Starting with the first term of \cref{eqn:terms12},
\begin{equation*}
\label{eqn:term1}
\begin{split}
&I\{ 0 \le A_0 < -(\widehat{A}_n - A_0), |\widehat{A}_n-A_0|\le \epsilon_n \}  + I\{ 0 \le A_0 < -(\widehat{A}_n - A_0), |\widehat{A}_n-A_0| > \epsilon_n \} \le \\&I\{ 0 \le A_0 < \epsilon_n\} + I \{ |\widehat{A}_n - A_0| > \epsilon_n \}, \\    
\end{split}
\end{equation*}
 where $I\{ 0 \le A_0 < \epsilon_n\} \rightarrow 0$ by Assumption \ifnum\jasa=1 10 \else \ref{assump:boundeddensity}\fi (bounded density assumption) and
$I \{ |\widehat{A}_n - A_0| > \epsilon_n \} \rightarrow 0$ by \cref{eqn:AnA0} (and definition of consistency). Specifically, assuming Assumption \ifnum\jasa=1 10\else \ref{assump:boundeddensity}\fi,
$Pr(0\le A_0 \le \epsilon_n) = \int^{\epsilon_n}_{0} f_{A_0}(a_0)da_0 \approx f_{a_0}(0)\epsilon_n \le b\epsilon_n \rightarrow 0 $ as $\epsilon_n \rightarrow 0 $ for some bound $b$. The second term can be shown to be bounded similarly. 
\begin{equation*}
\begin{split}
&I \{- (\widehat{A}_n - A_0) \le A_0 < 0 \} = \\
&I \{- (\widehat{A}_n - A_0) \le A_0 < 0 , |\widehat{A}_n - A_0| \le \epsilon_n\} +
I \{- (\widehat{A}_n - A_0) \le A_0 < 0 , |\widehat{A}_n - A_0| > \epsilon_n\} \leq\\
&I(-\epsilon_n \leq A_0 < 0) + I(|\widehat{A}_n - A_0| > \epsilon_n).
\end{split}
\end{equation*} The first term for the inequality is because $-(\widehat{A}_n - A_0) < 0$, $\widehat{A}_n - A_0$ must be positive; thus, $-\epsilon_n \leq \widehat{A}_n - A_0 \leq \epsilon_n$.  The second term is because the indicator is a subset.

Because of Equation \cref{eqn:AnA0}, $Pr(|\widehat{A}_n - A_0| > \epsilon_n) \rightarrow 0$, and $Pr(-\epsilon_n \leq A_0 < 0) \rightarrow 0$ due to Assumption \ifnum\jasa=1 10\else \ref{assump:boundeddensity}\fi.
Specifically, $I(-\epsilon_n \leq A_0 < 0)  \implies Pr(-\epsilon_n \le A_0 < 0) \approx \epsilon_n f_{a_0}(0) < \epsilon_n b \rightarrow 0$ as $\epsilon_n \rightarrow 0$.
This completes the proof of \cref{eqn:step1}.
\end{proof}
}

\begin{comment}
MARCH 2024 UPDATE: WE CAN IGNORE BELOW I THINK.

We want to use the below equation. I think he said that, as long as the density of the quantity $\valueoneTRUE - (1-\tol)\valueonemaxTRUE$ is bounded then the whole thing converges to 0?
\begin{equation}
%\begin{align}
    \lim_{n \to \infty} E\left(\left|I( \valueoneHAT - (1-\tol)\valueonemaxHAT \ge 0) - I( \valueoneTRUE - (1-\tol)\valueonemaxTRUE \ge 0 )\right|\right) = 0 \\
%\end{align}
\end{equation}
We want to show that $\valueoneHAT - (1-\tol)\valueonemaxHAT -\xrightarrow{p} \valueoneTRUE - (1-\tol)\valueonemaxTRUE$
Let $X_n \xrightarrow{p} \mu$ and $Y_n \xrightarrow{p} \gamma $. Then, $Pr(|X_n - \mu| > \epsilon \text{ or } |Y_n - \gamma|>\epsilon) \leq Pr(|X_n - \mu|>\epsilon) + Pr(|Y_n - \gamma|>\epsilon)$. If $n \rightarrow \infty$, then both $Pr(|X_n - \mu|>\epsilon) \rightarrow 0$ and $Pr(|Y_n - \gamma|>\epsilon) \rightarrow 0$. We want to use this as the generic idea for $X_n$ being OS and $Y_n$ being CIF.
\end{comment}

\begin{comment}
MARCH 2024 UPDATE: DEFINITELY IGNORE BELOW. PROOF IS COMPLETE.

\subsubsection{OLD STUFF: Notes from before 2/20/24 meeting}

Sample size $n$ such that $\hat{d}_n$. Let 
\[
V(\hat{d}_n) = \begin{pmatrix}
V_1(\hat{d}_n) \\ V_2(\hat{d}_n)
\end{pmatrix}
\]

We want to show:
\[
\lim_{{n \to \infty}} \bigg( \frac{V_1(d_*) - V_1(\hat{d}_n)}{V_1(d_*)} - \tol \bigg)_+ \xrightarrow{p} 0
\]
\\

Our optimal $d_{**}$ satisfies the following two criteria:
\begin{enumerate}
    \item $E\left[ I(\frac{V_1(d_*|Z)-V_1(d_{**}|Z)}{V_1(d_*|Z)} \leq \tol ) \right] = 1$
    \item maximize $E\left[V_2(d_{**}|Z) \;\middle|\; I(\frac{V_1(d_*|Z)-V_1(d_{**}|Z)}{V_1(d_*|Z)} \leq \tol ) = 1\right]$
\end{enumerate}
\end{comment}